%%% Основные сведения %%%
\newcommand{\thesisAuthorLastName}{Белокопытов}
\newcommand{\thesisAuthorOtherNames}{Антон Сергеевич}
\newcommand{\thesisAuthorInitials}{А.\,С.}
\newcommand{\thesisAuthor}             % Диссертация, ФИО автора
{%
	\texorpdfstring{%
		\thesisAuthorLastName~\thesisAuthorOtherNames%
	}{%
		\thesisAuthorLastName, \thesisAuthorOtherNames%
	}
}
\newcommand{\thesisAuthorShort}        % Диссертация, ФИО автора инициалами
{\thesisAuthorInitials~\thesisAuthorLastName}

%\newcommand{\thesisUdk}                % Диссертация, УДК (если не требуется, оставляем закомментированным)
%{УДК 159.9}

\newcommand{\thesisTitle}              % Диссертация, название
{Разработка алгоритмов поиска поведенчески релевантных нейрональных источников в многоканальных нейрофизиологических данных с использованием методов нелинейного вложения}

\newcommand{\thesisSpecialtyNumber}    % Диссертация, специальность, номер
{5.12.4} % Или укажите внутренний шифр ВШЭ, если он используется

\newcommand{\thesisSpecialtyTitle}     % Диссертация, специальность, название
{Когнитивное моделирование}

\newcommand{\thesisDegree}              % Диссертация, ученая степень
{кандидата когнитивных наук}

\newcommand{\thesisDegreeShort}         % Диссертация, ученая степень, краткая запись
{кандидат когнитивных наук}

\newcommand{\thesisCity}               % Диссертация, город написания диссертации
{Москва}

\newcommand{\thesisYear}               % Диссертация, год написания диссертации
{2029}

\newcommand{\thesisOrganization}       % Диссертация, организация
{Национальный исследовательский университет <<Высшая школа экономики>>}

\newcommand{\thesisOrganizationShort}  % Диссертация, краткое название организации для доклада
{НИУ ВШЭ}

\newcommand{\thesisInOrganization}     % Диссертация, организация в предложном падеже: Работа выполнена в ...
{Национальном исследовательском университете <<Высшая школа экономики>>}

%% Научный руководитель
\newcommand{\supervisorFio}              
{Осадчий Алексей Евгеньевич}
\newcommand{\supervisorRegalia}          
{доктор физико-математических наук}
\newcommand{\supervisorFioShort}         
{А.\,Е.~Осадчий}
\newcommand{\supervisorRegaliaShort}     
{д-р~физ.-мат.~наук}

%% Оппоненты / Комитет по защите (для ВШЭ) 
%% Оставьте заглушки до момента формирования комитета
\newcommand{\opponentOneFio}           
{Фамилия Имя Отчество}
\newcommand{\opponentOneRegalia}       
{доктор психологических наук, профессор}
\newcommand{\opponentOneJobPlace}      
{Название организации первого члена комитета}
\newcommand{\opponentOneJobPost}       
{главный научный сотрудник}

\newcommand{\opponentTwoFio}           
{Фамилия Имя Отчество}
\newcommand{\opponentTwoRegalia}       
{кандидат биологических наук}
\newcommand{\opponentTwoJobPlace}      
{Название организации второго члена комитета}
\newcommand{\opponentTwoJobPost}       
{старший научный сотрудник}

\newcommand{\leadingOrganizationTitle} 
{Название ведущей организации (если требуется)}

\newcommand{\defenseDate}              % Защита, дата
{DD месяца 202X~г.~в~XX:00}
\newcommand{\defenseCouncilNumber}     % Защита, номер диссертационного совета
{Диссертационный совет по когнитивным наукам}
\newcommand{\defenseCouncilTitle}      % Защита, учреждение диссертационного совета
{НИУ ВШЭ}
\newcommand{\defenseCouncilAddress}    % Защита, адрес учреждение диссертационного совета
{г. Москва, ул. Мясницкая, д. 20}
\newcommand{\defenseCouncilPhone}      % Телефон для справок
{+7~(495)~772-95-90}

\newcommand{\defenseSecretaryFio}      % Секретарь диссертационного совета, ФИО
{Фамилия Имя Отчество}
\newcommand{\defenseSecretaryRegalia}  % Секретарь диссертационного совета, регалии
{канд.~наук}            

\newcommand{\synopsisLibrary}          % Автореферат, название библиотеки
{Библиотека НИУ ВШЭ}
\newcommand{\synopsisDate}             % Автореферат, дата рассылки
{DD месяца 202X~года}

% Ключевые слова для метаданных PDF диссертации и автореферата
\providecommand{\keywords}
{электроэнцефалография, магнитоэнцефалография, линейная пространственная фильтрация, нелинейное вложение, UMAP, индуцированная активность, eSPoC, когнитивные науки}